\documentclass[11pt]{article}
\input{packages.tex}
\DeclareMathOperator{\sign}{sign}
\DeclareMathOperator{\Imm}{Im}

% Настройка колонтитулов
\pagestyle{fancy}
\fancyhf{} % Очистить все поля
\fancyhead[L]{Чил аут ребятки, копим свагу}
\fancyhead[R]{\href{https://t.me/ceagest}{Автор}}
\fancyfoot[C]{\thepage}               % Справа внизу

\setlength{\headheight}{13.6pt}

\hypersetup{
    colorlinks=true,
    linkcolor=blue,
    citecolor=blue,
    filecolor=blue,
    urlcolor=blue,
}

\begin{document}

\begin{titlepage}
    \centering
    \vspace*{\fill}  % заполняет пространство сверху

    {\Huge Лекция 1. Булева алгебра. КНФ и ДНФ. Полином Жегалкина. Классы Поста
    \par}

    \vspace*{\fill}  % заполняет пространство снизу
\end{titlepage}

\setcounter{page}{1}

\section{Булева алгебра. Формулы и отображения}

\subsection{Введение}

Приведём таблицу истинности для основных логических связок одной и двух переменных:

\begin{center}
    \begin{tabular}{|c|c|c|c|c|c|c|c|c|c|}
    \hline
    $x$ & $y$ & $\overline{x}$ & $x \land y$ & $x \lor y$ & $x \to y$ & $x \equiv y$ & $x \oplus y$ & $x \mid y$ & $x \downarrow y$ \\
    \hline
    0 & 0 & 1 & 0 & 0 & 1 & 1 & 0 & 1 & 1 \\
    \hline
    0 & 1 & 1 & 0 & 1 & 1 & 0 & 1 & 1 & 0 \\
    \hline
    1 & 0 & 0 & 0 & 1 & 0 & 0 & 1 & 1 & 0 \\
    \hline
    1 & 1 & 0 & 1 & 1 & 1 & 1 & 0 & 0 & 0 \\
    \hline
\end{tabular}
\end{center} \leavevmode
\\
Здесь $x \land y$ --- конъюнкция; $x \lor y$ --- дизъюнкция; $x \oplus y$ --- сложение по модулю $2$; $x \mid y$ --- штрих Шеффера; $x \downarrow y$ --- стрелка Пирса.
\\
В рамках этого раздела курса мы будем рассматривать булевы функции, то есть отображения 
из декартова произведения $\{0, 1\}^k$ (или, если угодно, из булева куба размерности $k$) в множество $\{0, 1\}$:
\[f \colon \{0, 1\}^k \to \{0, 1\}\]

\begin{lemma}[Законы де Моргана]
    \begin{equation*}
        \overline{\overline{x} \lor \overline{y}} = x \land y
        \quad \quad \quad 
        \overline{\overline{x} \land \overline{y}} = x \lor y
    \end{equation*}
\end{lemma}

\begin{proof}
    Выписываем таблицу истинности и получаем удовольствие.
\end{proof}

\begin{lemma}[Закон поглащения]
    \begin{equation*}
        x \lor (x \land y) = x
    \end{equation*}
\end{lemma}

\begin{proof}
    Ну здесь вообще очев, товарищи.
\end{proof}

\subsection{Немного о множествах}

\paragraph{Парадокс Рассела.}

Пусть $K$ --- множество всех множеств, не содержащих себя в качестве элемента. Содержит 
ли $K$ себя в качестве элемента?

\begin{figure}[h]
\centering
\includegraphics[width=0.55\linewidth]{pictures/Множества и индикаторные функции.jpg}
\end{figure} \leavevmode
\\
Пусть $U$ --- множество всех элементов, $A \subseteq U$, $B \subseteq U$. Пусть $a$ и $b$ --- индикаторные функции 
множеств $A$ и $B$ соответственно, то есть:
\begin{equation*}
    a(x) = 
    \begin{cases}
        1, & x \in A \\
        0, & x \not \in A
    \end{cases}
    \quad \quad \quad
    b(x) = 
    \begin{cases}
        1, & x \in B \\
        0, & x \not \in B
    \end{cases}
\end{equation*} \leavevmode
\\
Найдём все значения $x \in U$, при которых каждое из следующих выражений истинно:
\[\overline{a}(x) = 1 \iff x \in U \setminus A \quad \quad \quad a(x) \lor b(x) = 1 \iff x \in A \cup B\]
\[a(x) \land b(x) = 1 \iff x \in A \cap B \quad \quad \quad a(x) \to b(x) = 1 \iff \overline{A} \cup B\]

\section{Существенные и фиктивные переменные. КНФ и ДНФ}

\subsection{Существенные и фиктивные переменные. Литералы}

\begin{definition}
    Переменная $x_i$ функции $f(x_1, \ldots, x_n)$ называется \textbf{существенной},
    если существует набор значений переменных $\alpha_1, \ldots, \alpha_n$, где $\alpha_j \in \{0, 1\}$,
    такой, что:
    \[f\left(\alpha_1, \ldots, \alpha_{i - 1}, 0, \alpha_{i + 1}, \ldots, \alpha_n\right) \neq 
    f\left(\alpha_1, \ldots, \alpha_{i - 1}, 1, \alpha_{i + 1}, \ldots, \alpha_n\right)\]
    Если переменная не является существенной, то её называют \textbf{фиктивной}.
\end{definition}

\begin{definition}
    Переменная или её отрицание называется \textbf{литералом}. Дизъюнкция литералов называется
    \textbf{дизъюнктом}, конъюнкция --- \textbf{конъюнктом}.
\end{definition}

\subsection{КНФ и ДНФ. Их совершенные вариации}

Рассмотрим некоторую функцию $f$ трёх переменных и выпишем набор принимаемых ею значений для всевозможных наборов значений аргументов:

\begin{center}
    \begin{tabular}{|c|c|c|c|}
    \hline
    $x$ & $y$ & $z$ & $f$ \\
    \hline
    0 & 0 & 0 & 1 \\
    \hline
    0 & 0 & 1 & 0 \\
    \hline
    0 & 1 & 0 & 0 \\
    \hline
    0 & 1 & 1 & 1 \\
    \hline
    1 & 0 & 0 & 1 \\
    \hline
    1 & 0 & 1 & 0 \\
    \hline
    1 & 1 & 0 & 1 \\
    \hline
    1 & 1 & 1 & 1 \\
    \hline
\end{tabular}
\end{center} \leavevmode
\\
Мы хотим составить такую формулу, в которой использовались бы только связки $\land$, $\lor$, $\lnot$ и которая
задаёт функцию $f$. С этой целью рассмотрим все строки таблицы, в которых $f$ принимает значение $0$, и выпишем 
следующую конъюнкцию дизъюнктов по этим строкам:
\[(x \lor y \lor \overline{z}) \land (x \lor \overline{y} \lor z) \land (\overline{x} \lor y \lor \overline{z})\]
Несложно видеть, что каждая из скобок обнуляется тогда и только тогда, когда значения каждой из переменных
равны соответствующим значениям из строк с нулями $f$, по которым строился данных дизъюнкт. Это и означает, что 
полученная формула задаёт функцию $f$. При этом она называется \textbf{конъюнктивной нормальной формой (КНФ)}. Если 
в каждом дизъюнкте присутствуют все переменные (как и в рассматриваемом примере), то эта форма называется 
\textbf{совершенной (СКНФ)}. Ясно, что КНФ существует для всех функций, а СКНФ --- для всех функций, кроме тождественной единицы (так как строим по нулям).
\\
Теперь по аналогии рассмотрим все строки таблицы, в которых $f$ принимает значение $1$, и выпишем следующую 
дизъюнкцию конъюнктов по этим строкам:
\[(\overline{x} \land \overline{y} \land \overline{z}) \lor (\overline{x} \land y \land z) \lor (x \land \overline{y} \land \overline{z}) \lor (x \land y \land \overline{z}) \lor (x \land y \land z)\]
Несложно видеть, что каждая из скобок истинна тогда и только тогда, когда значения каждой из переменных
равны соответствующим значениям из строк с единицами $f$, по которым строился данных конъюнкт. Это и означает, что 
полученная формула задаёт функцию $f$. При этом она называется \textbf{дизъюнктивной нормальной формой (ДНФ)}. Если 
в каждом конъюнкте присутствуют все переменные (как и в рассматриваемом примере), то эта форма называется 
\textbf{совершенной (СДНФ)}. Ясно, что ДНФ существует для всех функций, а СДНФ --- для всех функций, кроме тождественного нуля (так как строим по единицам).

\section{Полином Жегалкина}

В данном параграфе для краткости будем заменять обозначения сложения по модулю $2$ ($\oplus$) и конъюнкции ($\land$) на 
$+$ и $\cdot$ соответственно. Также будем считать известными следующие свойства:
\[x \lor (y \land z) = (x \lor y) \land (x \lor z)\]
\[x \land (y \lor z) = (x \land y) \lor (x \land z)\]
\[x \land (y + z) = (x \land y) + (x \land z)\]
Заметим, что $\overline{x} = x + 1$. При этом в секции выше было показано, что любую булеву функцию 
можно выразить с помощью формулы, использующей лишь связки $\land$, $\lor$ и $\lnot$. Значит, пользуясь законами де Моргана, любую функцию 
можно выразить формулой, которая содержит лишь $+$, $\cdot$ и $1$.
\\
Пусть дана булева функция $f(x_1, \ldots, x_n)$. Построим для неё КНФ, заменим все отрицания на сложение по модулю $2$, а затем раскроем все скобки,
пользуясь вышеописанными свойствами (также учитываем, что $x \cdot x = x$, а значит у переменных не будет степеней, выше первой). В итоге получим формулу, 
называемую \textbf{полиномом Жегалкина} функции $f$:
\begin{equation*} \label{pol:1} \tag{$\ast$}
    \alpha_0 + \alpha_1 x_1 + \alpha_2 x_2 + \ldots + \alpha_n x_n + \alpha_{12} x_1 x_2 + \alpha_{13} x_1 x_3 + \ldots + \alpha_{123} x_1 x_2 x_3 + \ldots + \alpha_{123 \ldots n} x_1 x_2 x_3 \ldots x_n
\end{equation*}

\begin{remark}
    Важно понимать, что, если раскрыты не все скобки и/или не приведены подобные слагаемые, то получен не полином Жегалкина.
\end{remark}

\begin{theorem}
    Для каждой булевой функции существует единственный полином Жегалкина.
\end{theorem}

\begin{proof}
    Существование полинома следует из существования КНФ и рассуждений выше.
    Предположим от противного, что у некоторой функции $f(x_1, \ldots, x_n)$ существует два \textit{различных}
    полинома Жегалкина. Обозначим эти полиномы $g_f(x_1, \ldots, x_n)$ и $\widetilde{g}_f(x_1, \ldots, x_n)$.
    Будем почленно их сравнивать, начиная с одночлена $\alpha_0$ (здесь мы естественным образом полагаем, что полиномы 
    имеют именно вид (\ref{pol:1}), то есть слагаемые и переменные в произведениях упорядочены ровно таким образом).
    Пусть $\alpha x_{i_1} \ldots x_{i_k}$ и $\beta x_{i_1} \ldots x_{i_k}$ --- первые отличающиеся в этих полиномах для $f$ слагаемые (самые левые и наименьшей степени, где под степенью понимается количество входящих в произведение переменных).
    Так как сумма двух полиномов для $f$ есть тождественных ноль (так как при каждом фиксированном наборе переменных они принимают равные значения), а $\alpha + \beta = 1$ (в силу различности слагаемых, коэффициенты пред ними должны быть различны),
    то сумму полиномов можно выразить двумя способами:
    \begin{equation*} \label{eq:2} \tag{$\star$}
        0 = g_f(x_1, \ldots, x_n) + \widetilde{g}_f(x_1, \ldots, x_n) = x_{i_1} x_{i_2} \ldots x_{i_k} + \underbrace{\ldots \ldots \ldots \ldots \ldots}_{\text{степень не меньше } k}
    \end{equation*}
    Заметим, что полученные выше равенства выполнены вообще для любых наборов значений переменных $x_1, \ldots, x_n$.
    Так давайте же рассмотрим (\ref{eq:2}) для такого набора значений переменных, что $x_{i_1} = x_{i_2} = \ldots = x_{i_k} = 1$, а все остальные нули.
    Тогда выражение $x_{i_1} \ldots x_{i_k} = 1$, а каждое слагаемое из многоточия имеет степень не меньше $k$, и, следовательно, содержит в своём произведении хотя бы один нулевой сомножитель.
    Это означает, что всё выражение под многоточием обращается в ноль, а значит правая часть равна $1$, то есть для рассматриваемого набора значений переменных выполнено:
    \[0 = g_f(x_1, \ldots, x_n) + \widetilde{g}_f(x_1, \ldots, x_n) = 1 + 0 = 1\]
    Отсюда получаем противоречие, что завершает доказательство теоремы.
\end{proof}

\section{Полнота. Классы Поста}

\begin{definition}
    \textbf{Замыканием} множества функций $\Phi = \left\{f_1, \ldots, f_k\right\}$ называют 
    множество всех функций, которые можно получить из $\Phi$ при помощи суперпозиции и замены переменных, 
    в том числе с отождествлением. Говоря проще, можно $f(x, y)$ заменить на $f(x, z)$ или $f(x, x)$.
\end{definition}

\begin{definition}
    Множество функций называется \textbf{замкнутым}, если оно совпадает со своим замыканием.
\end{definition}

\begin{definition}
    Множество функций называют \textbf{полным}, если его замыкание совпадает с множеством всех булевых функций.
\end{definition}

\subsection{Классы Поста}

\paragraph{Функции, сохраняющие ноль.}
\[T_0 = \{f(x_1, \ldots, x_n) \mid f(0, \ldots, 0) = 0\}\]
\paragraph{Функции, сохраняющие единицу.}
\[T_1 = \{f(x_1, \ldots, x_n) \mid f(1, \ldots, 1) = 1\}\]
\paragraph{Линейные функции.}
\[L \text{ --- множество всех таких функций } f(x_1, \ldots, x_n) \text{, что для них}\]
\[\text{полином Жегалкина имеет не более, чем первую степень}\]
\paragraph{Самодвойственные функции.}
\[S = \{f(x_1, \ldots, x_n) \mid f(\overline{x_1}, \ldots, \overline{x_n}) = \overline{f(x_1, \ldots, x_n)}\}\]
\paragraph{Монотонные функции.}
\[M = \{f(x_1, \ldots, x_n) \mid \alpha \succeq \beta \implies f(\alpha) > f(\beta)\}\]
Здесь $"\succeq"$ --- частичный порядок, заданный на множестве всех $n$-ок из нулей и единиц. При этом 
\[\alpha \succeq \beta \iff \alpha_i \geq \beta_i \quad \forall i \in \{1, \ldots, n\}\]

\begin{theorem}[Поста]
    Множество функций полно тогда и только тогда, когда оно не содержится целиком ни в одном
    из классов Поста.
\end{theorem}

\end{document}